%%%% Modèle proposé par frederic.mazaleyrat@ens-paris-saclay.fr %%%%
%%%% 3 avril 2021 %%%%

\documentclass[french,12pt,a4paper]{book}
\usepackage[utf8]{inputenc}
\usepackage[T1]{fontenc}
\usepackage[french]{babel}
\usepackage{amsmath}
\usepackage{amsfonts}
\usepackage{fancyhdr}
\usepackage{amssymb}
\usepackage{color} % où xcolor selon l'installation
\definecolor{Prune}{RGB}{99,0,60}
\usepackage{mdframed}
\usepackage{multirow} %% Pour mettre un texte sur plusieurs rangées
\usepackage{multicol} %% Pour mettre un texte sur plusieurs colonnes
\usepackage{scrextend} %Forcer la 4eme  de couverture en page pair
\usepackage{tikz}
\usepackage{graphicx}
\usepackage[absolute]{textpos} 
\usepackage{colortbl}
\usepackage{array}
\usepackage{geometry}
\usepackage{hyperref}

\usepackage{lipsum} % à retirer!!!
\begin{document}

\begin{titlepage}


%\thispagestyle{empty}

\newgeometry{left=6cm,bottom=2cm, top=1cm, right=1cm}

\tikz[remember picture,overlay] \node[opacity=1,inner sep=0pt] at (-13mm,-135mm){\includegraphics{Bandeau_UPaS.pdf}};

% fonte sans empattement pour la page de titre
\fontfamily{fvs}\fontseries{m}\selectfont


%*****************************************************
%******** NUMÉRO D'ORDRE DE LA THÈSE À COMPLÉTER *****
%******** POUR LE SECOND DÉPOT                   *****
%*****************************************************

\color{white}

\begin{picture}(0,0)

\put(-110,-743){\rotatebox{90}{NNT: 2020UPASA000}}
\end{picture}
 
%*****************************************************
%**  LOGO  ÉTABLISSEMENT PARTENAIRE SI COTUTELLE
%**  CHANGER L'IMAGE PAR DÉFAUT **
%*****************************************************
\vspace{-10mm} % à ajuster en fonction de la hauteur du logo
\flushright \includegraphics[scale=1]{logo2.png}




%*****************************************************
%******************** TITRE **************************
%*****************************************************
\flushright
\vspace{10mm} % à régler éventuellement
\color{Prune}
\fontfamily{cmss}\fontseries{m}\fontsize{22}{26}\selectfont
  Titre de la thèse en français (sur plusieurs lignes si nécessaire)

\normalsize
\color{black}
\Large{\textit{Titre de la thèse en anglais}}
%*****************************************************

%\fontfamily{fvs}\fontseries{m}\fontsize{8}{12}\selectfont

\vspace{1.5cm}
\normalsize

\textbf{Thèse de doctorat de l'Université Paris-Saclay}


\vspace{15mm}

École doctorale n$^{\circ}$ 000, dénomination et sigle\\
\small Spécialité de doctorat: voir annexe\\
\footnotesize Unité de recherche: voir annexe\\
\footnotesize Référent: : voir annexe
\vspace{15mm}

\textbf{Thèse présentée et soutenue à .....,\\ le .... 202X, par}\\
\bigskip
\Large {\color{Prune} \textbf{Prénom NOM}}


%************************************
\vspace{\fill} % ALIGNER LE TABLEAU EN BAS DE PAGE
%************************************

%\flushleft \small \textbf{Composition du jury:}
\bigskip

\flushleft

\scriptsize
\begin{tabular}{|p{7cm}l}
\arrayrulecolor{Prune}
{\footnotesize \textbf{Composition du jury}}\\
& \\
\textbf{Prénom Nom} &   Président/e\\ 
Titre, Affiliation & \\
\textbf{Prénom Nom} &  Rapportrice \\ 
Titre, Affiliation   &   \\ 
\textbf{Prénom Nom} &  Rapporteur \\ 
Titre, Affiliation  &   \\ 
\textbf{Prénom Nom} &  Examinatrice \\ 
Titre, Affiliation   &   \\ 
\textbf{Prénom Nom} &  Examinateur \\ 
Titre, Affiliation   &   \\ 
\textbf{Prénom Nom} &  Examinateur \\ 
Titre, Affiliationt   &   \\ 

\end{tabular} 

\medskip
\begin{tabular}{|p{7cm}l}
\arrayrulecolor{Prune}
{\footnotesize \textbf{Direction de la thèse}}\\
& \\
\textbf{Prénom Nom} &   Directrice\\ 
Titre, Affiliation & \\
\textbf{Prénom Nom} &   Codirecteur\\ 
Titre, Affiliation  &   \\ 
\textbf{Prénom Nom} &   Coencadrante\\ 
Titre, Affiliation  &   \\ 
\textbf{Prénom Nom} &  tuteur en entreprise \\ 
Titre, Affiliation  &   \\ 


\end{tabular} 


\end{titlepage}

%%%%%%%%%%%%%%%%%%%%%%%%%%%%%%%%%%%%%%%%%%%%%%%%%%%%%%%%%%%%%%
% insérez ici les fichiers des chapitre

% \include{chapitre1}

\chapter{Quelques Conseils}
\section{Page de garde}
Ne changez rien dans le formatage, la taille des fontes, la couleur, l'alignement etc.

Les annexes, liste des ED, liste des référents etc, sont disponibles dans le modèle .docx sur le site de l'Université:

\url{https://www.universite-paris-saclay.fr/research/textes-de-reference/documents-de-reference-relatifs-la-soutenance-de-la-these#modelcover}

\section{Quatrième de couverture}

Les logos des ED sont à extraire du modèle .docx, même \verb!url!.

Ne pas réduire la police des résumés. Si ça ne rentre pas dans le cadre, résumez plus!

%%%%%%%%%%%%%%%%%%%%%%%%%%%%%%%%%%%%%%%%%%%%%%%%%%%%%%%%%%%%%%%
% 4eme de couverture
\ifthispageodd{\newpage\thispagestyle{empty}\null\newpage}{}
\thispagestyle{empty}
\newgeometry{top=1.5cm, bottom=1.25cm, left=2cm, right=2cm}
\fontfamily{rm}\selectfont

\lhead{}
\rhead{}
\rfoot{}
\cfoot{}
\lfoot{}

\noindent 
%*****************************************************
%***** LOGO DE L'ED À CHANGER IMPÉRATIVEMENT *********
%*****************************************************
\includegraphics[height=2.45cm]{EOBE}
\vspace{1cm}
%*****************************************************
\fontfamily{cmss}\fontseries{m}\selectfont

\small

\begin{mdframed}[linecolor=Prune,linewidth=1]

\textbf{Titre:} Titre de la thèse en français


\noindent \textbf{Mots clés:} Quelques mots-clé

\vspace{-.5cm}
\begin{multicols}{2}
\noindent \textbf{Résumé:}\lipsum[1-2] 
\end{multicols}

\end{mdframed}

\begin{mdframed}[linecolor=Prune,linewidth=1]

\textbf{Title:} Thesis title in English

\noindent \textbf{Keywords:} Some keywords

\begin{multicols}{2}
\noindent \textbf{Abstract:} \lipsum[1-2]
\end{multicols}
\end{mdframed}

%************************************
\vspace{\fill} % ALIGNER EN BAS DE PAGE
%************************************

\noindent
\color{Prune} \footnotesize Maison du doctorat de Université Paris-Saclay\\
2$^{\mathrm{e}}$ étage, aile ouest, École normale supérieure Paris-Saclay\\
4 avenue des Sciencs\\
91190 Gif-sur-Yvette, France

\end{document}